% ==================================================================
%  Section 2.6 – Cost Modeling
% ==================================================================
\section{Cost Modeling}
\label{sec:bg-cost-model}

Sections~\ref{sec:bg-mapping} and~\ref{sec:bg-mapspace} established that the
mapping problem requires choosing one mapping out of a combinatorially large
space and that the optimal choice depends on the \emph{cost} of each mapping (in terms of Energy-Delay-Product, Energy or Latency.
This section formalizes how that cost is estimated.

A \textbf{cost model} is a function that takes an architecture description and
a mapping and returns the cost without actually fabricating the chip or
running the workload on real hardware.
Two fundamentally different strategies exist.

\begin{itemize}
  \item \textbf{Simulation-based} models (e.g.\ cycle-accurate simulators)
        execute the data movement and compute operations on a software model of
        the hardware, tracking every memory access and pipeline event.
        They can be highly accurate, but their runtime grows with the number
        of operations in the workload, making them impractical when millions
        of mappings need to be evaluated.
  \item \textbf{Analytical} models compute the cost in closed form (or
        near--closed form) from the mapping's tiling factors, loop ordering,
        and the architecture's parameters.
        Their runtime is independent of the workload size, enabling rapid
        exploration of large map spaces.
        The trade-off is that simplifications, such as assuming perfect
        double-buffering or ignoring interconnect congestion, introduce
        approximation error.
\end{itemize}

Because map-space exploration (Section~\ref{sec:bg-mapspace}) requires
evaluating a cost function for a very large number of candidate mappings,
analytical models are the predominant choice in modern DNN mapping
frameworks~\cite{parashar2019timeloop, wu2019accelergy, kwon2019MAESTRO}.
The remainder of this section describes the three building blocks of such
models: the \emph{memory-access count}, the \emph{energy model}, and the
\emph{latency model}, together with the compound metrics derived from them.

% ------------------------------------------------------------------
\subsection{Memory-Operation Count}
\label{sec:mop-count}

The starting point of any analytical cost model is to determine how many
\emph{memory operations} (MOPs) each level of the hierarchy performs for a
given mapping.
A memory operation is a single scalar read or write from/to a storage level;
the total count depends on the tiling, loop ordering, and bypass
configuration of the mapping (Sections~\ref{sec:tiling}--\ref{sec:loop-ordering}).

Consider a memory level~$\ell$ that stores operands $\mathcal{O} \subseteq
\{\text{in}, \text{w}, \text{out}\}$ (the remaining operands bypass the
level).
At each iteration of the loops mapped on~$\ell$, a \emph{tile} of every
stored operand must be read from~$\ell$ and supplied to the levels below;
partial results are written back from below.
The total number of \emph{reads} and \emph{writes} at~$\ell$ can be
expressed as:
\begin{equation}
\label{eq:mops}
\text{reads}_\ell \;=\; \sum_{o \in \mathcal{O}}
  \frac{\text{tile}_{o}}{\text{reuse}_{o,\ell}}
  \;\cdot\; \prod_{k \in \mathcal{L}_\ell} f_k,
\qquad
\text{writes}_\ell \;=\; \text{fills}_\ell + \text{updates}_\ell,
\end{equation}
where $f_k$ are the tiling factors on level~$\ell$, the product ranges over
the loops~$\mathcal{L}_\ell$ mapped there, and $\text{reuse}_{o,\ell}$
captures temporal reuse from stationarity and halo overlap
(Section~\ref{sec:loop-ordering}).
The writes are the sum of \emph{fills}., data moved into the level from above level, and \emph{updates}, partial
sums drained from below levels, reads and writes are tracked separately so that read and write
bandwidths can be individually bounded.

Additionally, MOPs occurring at a fanout level (Section~\ref{sec:sa-abstraction})
are accounted for through \emph{multicast} and \emph{reduction}:
when an operand is spatially reused across PEs, a single read from the
level above can serve multiple instances (multicast), and corresponding
partial sums from multiple PEs may be aggregated before being written back
(reduction).

% ------------------------------------------------------------------
\subsection{Energy Model}
\label{sec:energy-model}

Given the MOP counts, the total energy is obtained by weighting each
access with the per-access energy of the corresponding storage level:
\begin{equation}
\label{eq:energy}
E \;=\;
  \underbrace{\sum_{\ell}\bigl(
    e_{\ell}^{\text{rd}}\cdot\text{reads}_\ell \;+\;
    e_{\ell}^{\text{wr}}\cdot\text{writes}_\ell
  \bigr)}_{\text{data-movement energy}}
  \;+\;
  \underbrace{\sum_{\ell} e_{\ell}^{\text{leak}}\cdot L}_{\text{leakage}}
  \;+\;
  \underbrace{e^{\text{MAC}}\cdot\text{MACs}}_{\text{compute energy}},
\end{equation}
where $e_{\ell}^{\text{rd}}$ and $e_{\ell}^{\text{wr}}$ are the per-scalar
read and write energies at level~$\ell$ (in pJ), $e_{\ell}^{\text{leak}}$
is the leakage power per clock cycle, $L$ is the total execution latency,
and $e^{\text{MAC}}$ is the energy of a single multiply-accumulate operation.
The summation runs over all memory levels and the compute level.

The per-access energies are technology-dependent constants that characterize
the hardware.
They can be provided directly for each level (e.g.\ measured or taken from a
data-sheet) or estimated by an external tool like \textbf{accelergy}~\cite{wu2019accelergy} a widely used energy-estimation
framework that queries technology-specific plug-ins, such as CACTI for
SRAMs~\cite{cacti}, to estimate the per-access energy and area of
each component given its configuration (depth, width, technology node, number
of ports, number of banks).
Because the per-access energies are constants of the architecture and do
\emph{not} change across mappings, they can be queried once during
architecture setup and then reused throughout exploration.
Compute energy is fixed once the workload is determined: the
total number of MACs is a property of the computation, not of how it is
mapped.
Since DRAM accesses are orders of magnitude more energy-expensive than
register-file accesses or multiply accumulations (${\sim}50/100\,\text{pJ}$ vs.\
${\sim}1\,\text{pJ}$ vs.\ ${\sim}0.2\,\text{pJ}$ per byte accessed at $45\,$nm%
~\cite{horowitz20141}), data movement dominates the energy budget, and so
minimizing the number of MOPs at the costliest levels is the primary
objective of any mapping optimization.

% ------------------------------------------------------------------
\subsection{Latency Model}
\label{sec:latency-model}

The latency of a mapping is the total number of clock cycles from the first
input read to the last output drain.
In an analytical model, latency is driven by two bottlenecks at every memory
level: the \emph{compute bound} and the \emph{bandwidth bound}.

\paragraph{Compute bound.}
If the available bandwidth is sufficient to keep the compute units fed, the
latency at level~$\ell$ equals the number of tiles processed times the
latency of one single tile:
\begin{equation}
\label{eq:lat-compute}
L_\ell^{\text{compute}} \;=\; c_{\text{tile}} \;\cdot\;
  \prod_{k \in \mathcal{L}_\ell} f_k,
\end{equation}
where $c_{\text{tile}}$ is the clock cycles required to process one tile
(propagated upward from the compute level).

\paragraph{Bandwidth bound.}
If the aggregate data rate demanded by a level exceeds its read or write
bandwidth~$B_\ell$, the level becomes a bottleneck and latency increases:
\begin{equation}
\label{eq:lat-bw}
L_\ell^{\text{BW,rd}}
  \;=\; \frac{\text{reads}_\ell + \text{drains}_\ell}{B_\ell^{\text{rd}}},
\qquad
L_\ell^{\text{BW,wr}}
  \;=\; \frac{\text{fills}_\ell + \text{updates}_\ell}{B_\ell^{\text{wr}}}.
\end{equation}
The effective latency contributed by each level is then
\begin{equation}
\label{eq:lat-level}
L_\ell \;=\; \max\!\bigl(L_\ell^{\text{compute}},\;
  L_\ell^{\text{BW,rd}},\; L_\ell^{\text{BW,wr}}\bigr),
\end{equation}
and the overall execution latency is the maximum across all levels:
\begin{equation}
\label{eq:latency}
L \;=\; \max_\ell\; L_\ell.
\end{equation}

When a level's bandwidth is saturated, the excess demand translates into
\emph{stall cycles} that idle the compute units:
$\Delta_\ell = L_\ell - L_\ell^{\text{compute}}$.
A mapping with zero stall cycles at every level achieves full
\emph{compute utilization}; in practice, the slowest memory level dictates
the overall throughput.

% ------------------------------------------------------------------
\subsection{Compound Metrics}
\label{sec:compound-metrics}

Energy and latency capture complementary aspects of a mapping's quality:
a mapping may reduce energy by reusing data more but increase latency due
to bandwidth stalls, or vice versa.
To compare mappings along both axes simultaneously, compound metrics are
commonly used.

\paragraph{Energy-Delay Product (EDP).}
The most widely adopted metric is the \emph{energy-delay product}:
\begin{equation}
\label{eq:edp}
\text{EDP} \;=\; E \times L \qquad [\text{pJ}\cdot\text{cc}].
\end{equation}
Minimizing the~EDP balances energy against latency: improving either
dimension while not degrading the other yields a better mapping.
EDP is used as the default optimization target by several frameworks,
including Timeloop~\cite{parashar2019timeloop}, and Factorflow \cite{Factorflow} .
